\documentclass{beamer}
\usepackage[french]{babel}
\usepackage[utf8]{inputenc}
\usepackage[T1]{fontenc}
\usepackage{graphicx}
\usepackage{tikz}
\usepackage{notes}

\usetheme{Madrid}
\usecolortheme{default}

\title{Mon Stage chez Allaw}
\subtitle{Une Aventure Humaine dans la Legal Tech}
\author{Malo Legendre-Lemaire}
\date{\today}
\institute{EPITA - Promo 2026}

\begin{document}

\frame{\titlepage}

% Slide 1: Qui est Allaw ?
\begin{frame}
\frametitle{Allaw}
\begin{center}
\includegraphics[width=0.6\textwidth]{../media/allaw.png}

\vspace{1cm}
{\Large Le Doctolib du Droit}
\end{center}

\note{
Bonjour, je m'appelle Malo et je vais vous présenter mon stage chez Allaw.

Allaw, c'est une jeune startup nantaise créée il y a un an par un père et son fils, Ludovic et Germain Stang. Leur ambition est simple mais ambitieuse : créer le "Doctolib du droit".

Pour ceux qui ne connaissent pas Doctolib, c'est la plateforme qui permet de prendre rendez-vous facilement avec des médecins. Allaw veut faire la même chose mais pour les professionnels du droit : avocats, notaires, commissaires de justice.

L'idée c'est de simplifier la vie des particuliers qui ont besoin d'un avocat, mais aussi celle des professionnels qui peuvent mieux gérer leur planning et leurs clients.
}
\end{frame}

% Slide 2: Comment j'ai découvert Allaw
\begin{frame}
\frametitle{La Rencontre}
\begin{center}
{\Large Février 2024}

\vspace{1cm}
Une rencontre grâce à des contacts communs

\vspace{1cm}
↓

\vspace{1cm}
Une journée d'immersion
\end{center}

\note{
Mon histoire avec Allaw a commencé en février 2024. Grâce à des contacts communs, j'ai eu la chance de rencontrer Ludovic Stang, le CEO d'Allaw.

Cette première rencontre s'est très bien passée et a rapidement débouché sur une journée d'immersion dans l'entreprise. C'était une occasion parfaite pour découvrir concrètement ce qui m'attendrait durant mon stage.

J'ai pu rencontrer une partie de l'équipe, qui était alors plus petite qu'aujourd'hui - ils étaient environ 6-7 personnes à l'époque. L'ambiance m'a tout de suite plu : jeune, dynamique, avec une vraie passion pour ce qu'ils font.
}
\end{frame}

% Slide 3: Pourquoi Allaw ?
\begin{frame}
\frametitle{Pourquoi Allaw ?}
\begin{itemize}
\item Une petite équipe en croissance (10 personnes)
\item Technologies que je maîtrise déjà
\item Impact direct sur le produit
\item Secteur en transformation
\item Nantes = chez moi
\end{itemize}

\note{
Plusieurs raisons m'ont poussé à choisir Allaw pour mon stage :

D'abord, le fait que ce soit une petite entreprise. Avec seulement 10 personnes, j'allais pouvoir découvrir le travail en équipe de façon très immersive. Chaque personne compte vraiment, et j'allais avoir un impact significatif sur le développement du produit.

Ensuite, techniquement, ils utilisaient des technologies que je maîtrisais déjà : Node.js, TypeScript, MongoDB. Ça me permettait de me concentrer sur la découverte du monde de l'entreprise plutôt que sur l'apprentissage technique.

Le secteur m'intéressait aussi beaucoup - la Legal Tech est en pleine transformation numérique, avec plein d'opportunités d'innovation.

Et puis, avantage non négligeable, l'entreprise est basée à Nantes et je suis nantais !

Enfin, j'étais vraiment séduit par leur mission : démocratiser l'accès au droit et simplifier la vie des professionnels juridiques. Ça correspond à mes valeurs.
}
\end{frame}

% Slide 4: L'équipe
\begin{frame}
\frametitle{L'Équipe Allaw}
\begin{center}
{\Large Une équipe passionnée}

\vspace{1cm}
CEO + CTO + Head of Sales

\vspace{0.5cm}
↓

\vspace{0.5cm}
Développeurs • Sales • Support • Communication
\end{center}

\note{
L'équipe d'Allaw, c'est vraiment ce qui fait la force de cette entreprise. 

Au sommet, on a Ludovic le CEO et son fils Germain qui est CTO, plus Cédric qui s'occupe des ventes.

En dessous, il y a nous, les développeurs - Nicolas en backend, Vincent et Titouan en frontend, plus moi comme stagiaire backend. Il y a aussi Yanis et Sarah qui s'occupent du business développement et du support client, et Nawel qui gère la communication.

Ce qui m'a frappé dès le début, c'est la passion de chacun pour ce qu'il fait. Tout le monde croit vraiment au projet et s'investit à fond. L'ambiance est super collaborative - il n'y a pas de hiérarchie pesante, on peut tous échanger librement.

J'ai aussi apprécié les stratégies mises en place pour la cohésion d'équipe : des sorties, des jeux, des moments de partage. Ça crée une vraie dynamique de groupe.
}
\end{frame}

% Slide 5: Mon intégration
\begin{frame}
\frametitle{Mon Intégration}
\begin{center}
{\Large Rapidement dans le bain}

\vspace{1cm}
🤝 Équipe bienveillante

\vspace{0.5cm}
🎯 Missions concrètes dès le début

\vspace{0.5cm}
💬 Échanges avec tous les services
\end{center}

\note{
Mon intégration s'est faite très naturellement. L'équipe est vraiment bienveillante et m'a fait sentir comme un membre à part entière dès le premier jour.

Ce qui m'a marqué, c'est qu'on m'a tout de suite confié des missions concrètes. Pas question de rester dans mon coin à faire de la veille technique - j'ai eu des vrais projets qui avaient un impact direct sur le produit.

J'ai aussi beaucoup apprécié de pouvoir échanger avec tous les services. Contrairement à une grosse entreprise où on reste souvent dans son silo, ici j'interagis régulièrement avec l'équipe sales, le support client, la communication. Ça me donne une vision globale de l'entreprise.

C'est d'ailleurs un des gros avantages d'être dans une petite structure : on comprend vraiment comment fonctionne une entreprise dans son ensemble.
}
\end{frame}

% Slide 6: Au cœur d'une création d'entreprise
\begin{frame}
\frametitle{Au Cœur d'une Création d'Entreprise}
\begin{center}
{\Large Découvrir toutes les étapes}

\vspace{1cm}
📈 Embauches

\vspace{0.5cm}
💰 Levées de fonds

\vspace{0.5cm}
🚀 Expansion
\end{center}

\note{
Ce qui rend cette expérience vraiment unique, c'est que je suis au cœur de la dynamique d'une création d'entreprise.

Pendant mon stage, j'ai pu voir toutes les étapes importantes pour l'expansion d'Allaw. Par exemple, j'ai assisté aux processus d'embauche - c'est fascinant de voir comment une équipe se construit, quels profils on recherche, comment on s'assure que les nouvelles personnes s'intègrent bien à la culture existante.

J'ai aussi été témoin des discussions autour des levées de fonds. C'est tout un monde que je ne connaissais pas - comprendre comment on présente son business model, comment on convainc des investisseurs, comment on planifie la croissance.

Et puis il y a toute la stratégie d'expansion : dans quelles villes se développer, comment approcher de nouveaux clients, comment adapter le produit selon les retours du terrain.

C'est vraiment formateur de voir tout ça de l'intérieur. Je contribue à la création des fondations techniques de l'entreprise, mais j'apprends aussi énormément sur l'aspect business.
}
\end{frame>

% Slide 7: Mes missions principales
\begin{frame}
\frametitle{Mes Missions}
\begin{center}
{\Large Développeur Backend}

\vspace{1cm}
🔔 Serveur de notifications

\vspace{0.5cm}
🌍 Gestion des décalages horaires

\vspace{0.5cm}
🔍 Barre de recherche globale

\vspace{0.5cm}
🔗 Maillage pour le SEO
\end{center}

\note{
Côté technique, mon rôle principal était d'aider Nicolas, l'unique développeur backend, dans ses tâches.

Mon projet principal, c'était la création d'un serveur de notifications. L'idée, c'est que quand quelqu'un prend rendez-vous sur le site, le professionnel reçoit une notification en temps réel sans avoir à actualiser sa page. C'était un vrai défi technique passionnant.

J'ai aussi travaillé sur la gestion des décalages horaires. Ça peut paraître simple, mais c'est crucial pour permettre à Allaw de s'étendre dans les DOM-TOM. Il faut que quelqu'un de Martinique puisse prendre rendez-vous avec un avocat de métropole sans problème de fuseaux horaires.

Autre projet : une barre de recherche globale qui permet de chercher naturellement. Par exemple, taper "rendez-vous avec Malo le 5 janvier" et que le système comprenne et trouve les bons résultats.

Et enfin, le maillage pour le SEO - connecter intelligemment les pages des professionnels pour améliorer le référencement Google.

Ce qui est important de noter, c'est que je me concentrais sur le backend, tandis que mes collègues s'occupaient du frontend. Vraie collaboration d'équipe.
}
\end{frame>

% Slide 8: Les défis humains
\begin{frame}
\frametitle{Les Défis Humains}
\begin{center}
{\Large Apprendre à travailler en équipe}

\vspace{1cm}
🗣️ Communication

\vspace{0.5cm}
🤝 Collaboration

\vspace{0.5cm}
⚖️ Équilibre autonomie/guidance
\end{center}

\note{
Au-delà des aspects techniques, j'ai dû apprendre à travailler vraiment en équipe, ce qui était nouveau pour moi.

La communication, c'est quelque chose qu'on sous-estime souvent quand on est étudiant. Ici, il faut savoir expliquer ce qu'on fait, demander de l'aide quand on bloque, tenir les autres au courant de ses avancées. J'ai appris à être plus clair dans mes explications et à poser les bonnes questions.

La collaboration aussi, c'est tout un art. Mon serveur de notifications devait s'interfacer avec le travail de Nicolas côté backend et celui de Vincent côté frontend. Il faut se coordonner, respecter les délais des uns et des autres, s'adapter aux contraintes de chacun.

Et puis il y a cet équilibre délicat entre autonomie et guidance. On m'a fait confiance sur des projets importants, mais je pouvais toujours compter sur l'aide de mes collègues quand j'en avais besoin. C'est formateur de se sentir responsable tout en sachant qu'on n'est pas seul.
}
\end{frame>

% Slide 9: Ce que j'ai appris sur moi
\begin{frame}
\frametitle{Découverte de Soi}
\begin{center}
{\Large J'ai gagné en...}

\vspace{1cm}
💪 Confiance en moi

\vspace{0.5cm}
🎯 Autonomie

\vspace{0.5cm}
👥 Capacité à travailler en équipe
\end{center}

\note{
Ce stage m'a beaucoup appris sur moi-même.

D'abord, j'ai gagné en confiance. Quand on me confie un projet comme le serveur de notifications et qu'on me fait confiance pour le mener à bout, ça vous change. Je me suis rendu compte que j'étais capable de gérer des projets complexes et de proposer des solutions techniques pertinentes.

J'ai aussi développé mon autonomie. Au début, j'avais tendance à demander de l'aide pour tout. Petit à petit, j'ai appris à chercher par moi-même, à tester différentes approches, à prendre des initiatives. Bien sûr, je demande encore conseil, mais je suis beaucoup plus débrouillard qu'avant.

Et puis ma capacité à travailler en équipe s'est vraiment développée. J'ai appris à écouter les autres, à intégrer leurs contraintes dans mes décisions, à communiquer efficacement. C'est quelque chose qu'on n'apprend pas vraiment à l'école, mais qui est essentiel dans le monde professionnel.
}
\end{frame>

% Slide 10: L'ambiance de travail
\begin{frame}
\frametitle{L'Ambiance de Travail}
\begin{center}
{\Large Une culture d'entreprise forte}

\vspace{1cm}
🎉 Jeune et dynamique

\vspace{0.5cm}
💡 Innovation et partage

\vspace{0.5cm}
🤗 Bienveillance
\end{center}

\note{
L'ambiance de travail chez Allaw, c'est vraiment quelque chose de spécial. La culture d'entreprise est très forte, axée sur l'innovation, le partage et l'humain.

C'est jeune et dynamique, sans être superficiel. On travaille dur, mais dans la bonne humeur. Il y a une vraie émulation collective - quand quelqu'un a une idée, on l'écoute, on la discute, on la teste.

L'innovation fait partie de l'ADN de l'entreprise. On nous encourage à proposer des améliorations, à tester de nouvelles approches. Mon serveur de notifications, par exemple, j'étais complètement libre sur les choix techniques. Cette liberté, c'est très motivant.

Et puis il y a cette bienveillance qui transpire partout. Les erreurs ne sont pas dramatisées, on apprend ensemble. Quand quelqu'un galère sur un problème, tout le monde se mobilise pour l'aider.

C'est exactement l'ambiance que j'espérais pour ma première expérience professionnelle. Ça correspond parfaitement à mes attentes et à ma vision du travail.
}
\end{frame}

% Slide 11: Les moments marquants
\begin{frame}
\frametitle{Moments Marquants}
\begin{center}
{\Large Des souvenirs impérissables}

\vspace{1cm}
🎮 Sorties d'équipe

\vspace{0.5cm}
🎯 Premiers projets aboutis

\vspace{0.5cm}
🤝 Relations créées
\end{center}

\note{
Il y a plein de moments qui resteront gravés dans ma mémoire.

Les sorties d'équipe, déjà. Des sessions de jeux, des afterworks, des moments plus informels où on apprend à se connaître autrement que par le travail. C'est là qu'on voit vraiment la personnalité de chacun et que se créent les vrais liens.

Mes premiers projets aboutis aussi. La première fois que mon serveur de notifications a fonctionné et qu'on a vu une notification s'afficher en temps réel, c'était un moment de fierté intense. Pareil quand j'ai résolu le problème des décalages horaires - voir que mon code était utilisé en production, ça n'a pas de prix.

Et puis toutes ces relations créées. Avec Nicolas qui m'a énormément appris côté technique, avec Vincent et Titouan avec qui j'ai découvert les joies du travail frontend-backend, avec toute l'équipe en fait. Ces personnes, je sais que je vais garder le contact avec elles.

C'est ça aussi la richesse d'un stage - les relations humaines qu'on construit et qui nous accompagnent longtemps après.
}
\end{frame>

% Slide 12: Impact sur mon parcours
\begin{frame}
\frametitle{Impact sur Mon Parcours}
\begin{center}
{\Large Cette expérience m'a...}

\vspace{1cm}
🎯 Confirmé mes choix d'orientation

\vspace{0.5cm}
🛠️ Donné de l'expérience concrète

\vspace{0.5cm}
🌱 Fait grandir humainement
\end{center}

\note{
Ce stage a eu un impact énorme sur mon parcours et ma vision de l'avenir.

D'abord, ça a confirmé mes choix d'orientation. J'aime vraiment le développement backend, j'aime résoudre des problèmes techniques complexes, j'aime voir mes solutions utilisées en production. Cette expérience m'a conforté dans l'idée que c'est ce que je veux faire plus tard.

Ensuite, ça m'a donné une expérience concrète inestimable. Maintenant, quand on me parle de travail en équipe, de gestion de projet, de contraintes de production, je sais de quoi il s'agit. J'ai vécu tout ça de l'intérieur.

Mais surtout, ça m'a fait grandir humainement. J'ai une meilleure connaissance de moi-même, de mes capacités, de mes points d'amélioration. J'ai appris à gérer le stress, à communiquer mieux, à être plus organisé.

Cette expérience va définitivement m'aider pour la suite de mes études et pour mes futurs stages ou emplois. J'ai maintenant une base solide sur laquelle construire.
}
\end{frame}

% Slide 13: Ce que je retiens
\begin{frame}
\frametitle{Ce que Je Retiens}
\begin{center}
{\Large L'humain avant tout}

\vspace{1cm}
❤️ Une équipe passionnée

\vspace{0.5cm}
🎯 Des projets avec du sens

\vspace{0.5cm}
🌟 Une expérience transformatrice
\end{center}

\note{
Si je devais retenir une seule chose de ce stage, ce serait l'importance de l'humain dans le travail.

Chez Allaw, j'ai découvert une équipe vraiment passionnée par ce qu'elle fait. Ce n'est pas juste un travail, c'est une mission commune : révolutionner l'accès au droit. Cette passion, elle est contagieuse et elle donne envie de se dépasser.

J'ai travaillé sur des projets qui ont du sens. Mon serveur de notifications, par exemple, ça va vraiment améliorer la vie des professionnels du droit au quotidien. Savoir que son travail a un impact concret, c'est très motivant.

Et puis cette expérience a été transformatrice pour moi. Je ne suis plus le même qu'à mon arrivée. J'ai grandi professionnellement et personnellement. J'ai découvert le plaisir de travailler dans une équipe soudée sur des projets ambitieux.

C'est exactement ce que j'espérais d'un premier stage : une vraie immersion dans le monde du travail, avec des responsabilités réelles et une équipe bienveillante pour m'accompagner.
}
\end{frame}

% Slide 14: Remerciements
\begin{frame}
\frametitle{Remerciements}
\begin{center}
{\Large Merci à toute l'équipe Allaw}

\vspace{1cm}
🙏 Pour votre accueil

\vspace{0.5cm}
🤝 Pour votre confiance

\vspace{0.5cm}
💪 Pour m'avoir fait grandir
\end{center}

\note{
Je voudrais terminer en remerciant chaleureusement toute l'équipe d'Allaw.

Merci pour votre accueil. Dès le premier jour, vous m'avez fait sentir comme un membre à part entière de l'équipe. Cette bienveillance, ce n'est pas donné partout.

Merci pour votre confiance. Me confier des projets importants comme le serveur de notifications, c'était un pari sur mes capacités. Cette confiance m'a poussé à me dépasser.

Et merci pour m'avoir fait grandir. Techniquement bien sûr, mais surtout humainement. J'ai appris énormément à vos côtés et je suis convaincu que ces apprentissages me seront précieux pour la suite.

Cette expérience chez Allaw restera gravée comme une aventure humaine exceptionnelle. Vous m'avez donné goût au travail en startup et confirmé que c'est dans ce type d'environnement que je veux évoluer.

Merci à tous !
}
\end{frame}

% Slide 15: Questions
\begin{frame}
\frametitle{}
\begin{center}
{\Huge Questions ?}

\vspace{2cm}
{\large Malo Legendre-Lemaire}\\
{\small malo.legendre-lemaire@epita.fr}
\end{center}

\note{
Merci pour votre attention. Je suis maintenant disponible pour répondre à vos questions, que ce soit sur mon expérience chez Allaw, les projets techniques sur lesquels j'ai travaillé, ou plus généralement sur la vie en startup.

N'hésitez pas, je serai ravi d'échanger avec vous !
}
\end{frame}

\end{document}